\documentclass[10pt]{beamer}

\usepackage[utf8]{inputenc}
\usepackage[spanish]{babel}
\usepackage{graphicx}

\mode<presentation>
\usetheme{Madrid}
%\usecolortheme[RGB={111,73,135}]{structure}
\usecolortheme[RGB={128,0,0}]{structure}
%\usecolortheme[RGB={0,96,0}]{structure}
%\usecolortheme[RGB={200,0,200}]{structure}
%\usecolortheme[RGB={0,128,0}]{structure}
%\usecolortheme[RGB={0,0,128}]{structure}
\usefonttheme{serif}
\useinnertheme{rectangles}
\useoutertheme{split}
\setbeamercovered{transparent}


% Definiciones para usar luego
\newtheorem{ejemplo}{Ejemplo}
\newtheorem{definicion}{Definición}

\title{Entendiendo el Kernel de Linux}
\author{Jesús Espino García\\Samuel Rodríguez Sevilla}
\date{13 de Marzo de 2009}
\subject{Entendiendo el Kernel de Linux}

\institute[GUL UC3M]{
  Grupo de Usuarios de Linux\\
  Universidad Carlos III de Madrid.\\
  \ \\
  \includegraphics[height=1.5cm]{imgs/gul}}

\setcounter{tocdepth}{2}

\AtBeginSubsection[]
{
  \begin{frame}<beamer>{Indice}
    \tableofcontents[sectionstyle=show/hide,subsectionstyle=show/shaded/hide]
  \end{frame}
}

\begin{document}

  \frame{\maketitle}
  \begin{frame}
	\setcounter{tocdepth}{1}
  	\tableofcontents[sectionstyle=show/show,subsectionstyle=show/shaded/hide]
	\setcounter{tocdepth}{2}
  \end{frame}
  
  \section{Introducción}
\subsection{¿Que es?}
\begin{frame}
	\begin{itemize}
		\item Núcleo del sistema operativo GNU/Linux.
		\item GPL.
		\item Estable.
		\item Maduro.
		\item Moderno.
		\item Eficiente.
		\item Versátil.
		\item Configurable.
		\item y mucho mas!!
	\end{itemize}
\end{frame}

\subsection{Un poco de historia}
\begin{frame}
	\begin{itemize}
		\item En Abril de 1991 Linus Torvals empieza el desarrollo de Linux.
		\item En Septiembre del 1991 lanza la versión 0.01
		\item En Diciembre del 1991 libera la primera versión ''self hosted''.
		\item En Marzo del 1994 aparece la versión 1.0.0.
		\item En 1996 aparece Tux como mascota de Linux y la versión 2.0 del mismo.
		\item En 1999 lanza la versión 2.2
		\item En 2001 lanza la versión 2.4
		\item En 2003 lanza la versión 2.6 (la actual).
	\end{itemize}
\end{frame}

\subsection{Numeración de versiones}
\begin{frame}
	\begin{block}{Hasta la versión 2.5}
		\begin{itemize}
			\item El primer numero indicaba la versión de la rama del kernel.
			\item El segundo indicaba si era estable o no (par estable, impar inestable).
			\item El tercero era la minor-release de esa rama.
		\end{itemize}
	\end{block}
	\begin{block}{versión 2.6.0 y posteriores}
		\begin{itemize}
			\item El primer y segundo numero se han estabilizado (no se esperan cambios a corto plazo).
			\item El tercer numero indica la versión del kernel.
			\item Si esta versión esta seguida por -rcX significa que es una versión de desarrollo.
			\item Si esta versión esta seguida por un .X significa que es estable y que esta en el bugfix X.
		\end{itemize}
	\end{block}
\end{frame}

\subsection{Git}
\begin{frame}{Git}
	\begin{itemize}
		\item Sistema de control de versiones del kernel de Linux.
		\item Distribuido.
		\item Desarrollado directamente por Linus Torvals.
		\item Apareció como una alternativa a BitKeeper.
		\item Existen muchas ramas desarrolladas de manera distribuida.
	\end{itemize}
\end{frame}

\subsection{Visión General}
\begin{frame}{A vista de pájaro}
	\begin{center}
		\includegraphics[height=6.5cm]{imgs/vista1}
	\end{center}
\end{frame}

\begin{frame}{Capas}
	\begin{center}
		\includegraphics[height=6.5cm]{imgs/vista2}
	\end{center}
\end{frame}

\begin{frame}{Subsistemas}
	\begin{center}
		\includegraphics[height=6.5cm]{imgs/vista3}
	\end{center}
\end{frame}

  \include{linux_mm}
  \include{linux_interrup}
  \include{linux_process}
  \include{linux_syscalls}
  \section{Input/Output}
\subsection{Dispositivos}
\begin{frame}{Dispositivos}
	\begin{itemize}
		\item Los dispositivos son todos los elementos de un ordenador que permiten realizar operaciones de entrada/salida. Por ejemplo los discos duros, el teclado, etc.
		\item Los dispositivos se representan en \emph{sysfs} con un sistema de \emph{kobjects} (se encuentra declarado en \emph{linux/kobject.h}). Estos \emph{kobjects} describen la información sobre un dispositivo para que puede ser fácilmente consultable.
		\item Para poder acceder a los dispositivos se necesitan los ficheros de dispositivo (típicamente se encuentran en \emph{/dev/}). Estos vienen descritos por dos números denominados \emph{major} y \emph{minor}  y por su tipo, que puede ser caracter o bloque.
	\end{itemize}
\end{frame}

\begin{frame}{Dispositivos}
	\begin{itemize}
		\item El \emph{major} identifica el tipo de dispositivo y el \emph{minor} identifica un dispositivo concreto de ese tipo.
		\item Los dispositivos permiten realizar al programador una serie de acciones (\emph{open}, \emph{read}, \emph{lseek}, \emph{ioctl}, etc.).
		\item Los controladores de dispositivo deben estar registrados en el sistema y utilizan un descriptor de tipo \emph{device\_driver} que debe encontrarse en dentro de una estructura de tipo \emph{device}.
		\item Los controladores de dispositivo compilados estáticamente en el núcleo se registran durante el proceso de inicio. El resto deberán llamar a la función de registro que le corresponda (por ejemplo \emph{pci\_register\_driver}).
		\item Casi todos los controladores de dispositivo siguen una política de asignación de recursos tardía. De este modo, asignan la memoria cuando se les solicita hacer algo (normalmente por una llamada a \emph{open}) y la liberan cuando ya no es necesaria (no hay ninguna aplicación utilizando sus recursos).
	\end{itemize}
\end{frame}

\subsection{Arquitectura de I/O en bloques}
\subsubsection{Arquitectura}
\begin{frame}{Arquitectura}
	Hay dos tipos de dispositivos:
	\begin{itemize}
		\item Dispositivos de caracter: leen y/o escriben un flujo de datos. Su acceso es secuencial. Ejemplos de este tipo de dispositivos son el ratón y el teclado.
		\item Dispositivos de bloque: leen y/o escriben bloques de datos y permiten el acceso aleatorio a los mismos. Ejemplos de este tipo de dispositivos tenemos los discos duros, DVDs, etc.
	\end{itemize}
\end{frame}

\begin{frame}{Arquitectura}
	Los dispositivos de bloques entrañan más complejidad.
	\begin{itemize}
		\item Los dispositivos se dividen en sectores.
		\item Los sectores se agrupan en bloques.
		\item Los bloques se agrupan en segmentos (un segmento es la cantidad de información que puede mandarse por un canal DMA).
		\item Los segmentos se agrupan en páginas (pueden tener bloques que no se encuentran en un segmento).
	\end{itemize}
\end{frame}

\begin{frame}{Arquitectura}
	\begin{itemize}
		\item Los bloques se leen a memoria en áreas denominadas buffers.
		\item Los buffers está definidos por la estructura \emph{buffer\_head} que se encuentra en \emph{linux/buffer\_head.h}.
		\item Esta estructura controla los datos sobre a qué bloque corresponde el buffer, de qué dispositivo, en qué estado se encuentra, cuantas referencias hay a él, etc.
	\end{itemize}
\end{frame}

\begin{frame}{Arquitectura}
	\begin{itemize}
		\item Las operaciones de I/O que se están realizando se encuentran definidas por la estructura \emph{bio} que se encuentra en \emph{linux/bio.h}.
		\item Esta estructura representa las operaciones sobre páginas.
		\item La operación de I/O puede afectar a más de una página. Se controla con el array \emph{bio\_vec} de \emph{bio}.
		\item Las operaciones se encolan un una cola de peticiones pendientes.
	\end{itemize}
\end{frame}

\subsubsection{IO Scheduler}
\begin{frame}{IO Scheduler}
	\begin{itemize}
		\item Las operaciones de I/O no son óptimas en tiempo.
		\item El sistema operativo planifica estas operaciones para mejorar el rendimiento.
		\item La planificación se lleva a cabo siguiendo alguna de las siguientes políticas (en el núcleo se denominan \emph{elevators}):
		\begin{itemize}
			\item \emph{Linus Elevator}
			\item \emph{Deadline I/O Scheduler}
			\item \emph{Anticipatory I/O Scheduler}
			\item \emph{Complete Fair Queuing I/O Scheduler}
			\item \emph{Noop I/O Scheduler}
		\end{itemize}
	\end{itemize}
\end{frame}

\subsection{VFS}

  \include{linux_vfs}
  \section{Inicio del sistema}
\subsection{Inicio del sistema}
\begin{frame}{Inicio del sistema}
	\begin{itemize}
		\item Es el proceso que permite poner en funcionamiento al sistema operativo.
		\item Comienza por la BIOS, que realiza una comprobación del hardware y pasa el control a un programa que típicamente se encuentra en el sector de inicio del disco duro.
		\item El cargador de arranque se encarga de cargar el núcleo en memoria y pasarle el control junto con una serie de parámetros (se llama a la función \emph{setup( )} que se encuentra desplazada 0x200 bytes del inicio del núcleo).
	\end{itemize}
\end{frame}

\subsection{Proceso de inicio}
\begin{frame}{Proceso de inicio}
	El proceso de arranque realiza las siguientes tareas:
	\begin{itemize}
		\item Se invoca a la BIOS para saber la cantidad de RAM del sistema o, en sistemas ACPI, obtener un mapa físico de la RAM.
		\item Ajusta la la velocidad del repetición del teclado.
		\item Inicializa el adaptador de vídeo.
		\item Reinicia el controlador de disco y obtiene los parámetros de los discos.
		\item Comprueba que haya un bus IBM Micro Channel.
		\item Comprueba si hay un ratón PS/2.
		\item Comprueba si la BIOS soporta APM (\emph{Advanced Power Management}).
		\item En caso de que la BIOS soporte EDD (\emph{Enhanced Disk Drive Services}), se solicita crear una tabla en memoria con la descripción de los discos duros.
	\end{itemize}
\end{frame}
\begin{frame}{Proceso de inicio}
	\begin{itemize}
 		\item Según si el núcleo era sin comprimir o comprimido, procede a continuar o desplaza el núcleo en memoria para colocarlo donde le corresponde.
		\item Se activa el pin A20 en el controlador de teclado 8042.
		\item Se crea una tabla de descriptores de interrupciones (IDT) provisional y una tabla global de descripción (GDT) provisional.
		\item Reinicia la unidad de coma flotante (FPU).
		\item Se reprograma el controlador programable de interrupciones (PIC) para enmascarar todas las interrupciones menos la de IRQ2.
		\item Se pone la CPU en modo protegido.
		\item Se llama a la función \emph{startup\_32( )}.
	\end{itemize}
\end{frame}

\subsection{startup\_32}
\begin{frame}{La función startup\_32 (I)}
	\begin{itemize}
		\item Inicializa los registros de segmentación y una pila provisional.
		\item Pone a 0 los bits del registro \emph{eflags}.
		\item Inicializa con ceros el área de datos del núcleo que se encuentran en \emph{\_edata} y \emph{\_end}.
		\item Llama a \emph{decompress\_kernel( )}.
		\item Una vez descomprimido el núcleo se llama a un nuevo método \emph{startup\_32~(~)}.
	\end{itemize}
\end{frame}
\begin{frame}{La función startup\_32 (y II)}
	El nuevo \emph{startup\_32( )} hace las siguiente operaciones:
	\begin{itemize}
		\item Inicializa los registros de segmentación con los valores finales.
		\item Rellena el segmento \emph{bss} del núcleo con ceros.
		\item Inicializa la tabla de páginas provisional del núcleo.
		\item Almacena la dirección del \emph{Page Global Directory} en el registro \emph{cr3}.
		\item Asigna la pila de modo núcleo al proceso 0.
		\item Vuelve a poner a 0 el registro \emph{eflags}.
		\item Llama a \emph{setup\_idt( )}.
		\item Almacena los parámetros obtenidos de la BIOS en el primer \emph{frame} de memoria.
		\item Identifica el modelo del procesador.
		\item Carga los registros de la tabla GDT e IDT.
		\item salta a la función \emph{start\_kernel( )}. 
	\end{itemize}
\end{frame}

\subsection{start\_kernel}
\begin{frame}{La función start\_kernel}
	Llama a las siguientes funciones:
	\begin{itemize}
		\item \emph{sched\_init( )}.
		\item \emph{build\_all\_zonelists( )}.
		\item \emph{page\_alloc\_init( )} y \emph{mem\_init( )}.
		\item \emph{trap\_init( )}.
		\item \emph{softirq\_init( )}.
		\item \emph{time\_init( )}.
		\item \emph{kmem\_cache\_init( )}.
		\item \emph{calibrate\_delay( )}.
		\item \emph{kernel\_thread( )}.
	\end{itemize}
\end{frame}


  \include{linux_fin}

\end{document}
